考虑 \(N\) 个周期边界 \(s_i=\pm1\) 的最近邻Ising自旋：\(H=-J\sum_i s_is_{i+1}-\mu B\sum_i s_i\)，其中 \(J,\mu,T>0\)。记 \(K=J/(k_BT)\)、\(h=\mu B/(k_BT)\)。
(1) 构建转移矩阵、求两个本征值与有限 \(N\) 的精确配分函数，并推出热力学极限自由能。
(2) 由最大本征值求磁化、零场磁化率，严格判定任意正温度是否有自发磁化。
(3) 以转移矩阵本征通道计算零场两点关联及关联长度，并给出低温和高温极限。
(4) 从同一配分函数求零场每格点能量和热容，用独立的涨落恒等式核对正定性并讨论极限顺序。
